\documentclass[11pt,a4paper]{article}
\usepackage[margin=25mm,headheight=14pt]{geometry}
\usepackage[T1]{fontenc}
\usepackage[utf8]{inputenc}
\usepackage{lmodern}
\usepackage{amsmath,amssymb,amsthm,mathtools}
\usepackage{microtype,booktabs,array,longtable,enumitem}
\usepackage{xcolor,hyperref,fancyhdr}
\definecolor{linkink}{HTML}{164F63}
\hypersetup{colorlinks=true,linkcolor=linkink,citecolor=linkink,urlcolor=linkink,
 pdftitle={Sparse Grid Moments and Quadratic Closure},
 pdfauthor={Research draft prepared with ChatGPT for the ProveIt project},
 pdfsubject={Exact finite-field laws, sharp obstructions, and cube rank enumerators}}
\setlength{\emergencystretch}{2em}
\setlength{\parskip}{3pt}
\setlist{itemsep=2pt,topsep=4pt}
\numberwithin{equation}{section}
\newtheorem{theorem}{Theorem}[section]
\newtheorem{lemma}[theorem]{Lemma}
\newtheorem{proposition}[theorem]{Proposition}
\newtheorem{corollary}[theorem]{Corollary}
\theoremstyle{definition}
\newtheorem{definition}[theorem]{Definition}
\newtheorem{example}[theorem]{Example}
\newtheorem{question}[theorem]{Question}
\theoremstyle{remark}
\newtheorem{remark}[theorem]{Remark}
\newcommand{\F}{\mathbb F}
\newcommand{\E}{\mathbb E}
\newcommand{\PP}{\mathbb P}
\newcommand{\1}{\mathbf 1}
\newcommand{\rank}{\operatorname{rank}}
\newcommand{\Span}{\operatorname{span}}
\newcommand{\supp}{\operatorname{supp}}
\newcommand{\cl}{\operatorname{cl}}
\newcommand{\TV}{\operatorname{TV}}
\newcommand{\Law}{\mathcal L}
\newcommand{\Mat}{\operatorname{Mat}}
\newcommand{\diag}{\operatorname{diag}}
\newcommand{\Tr}{\operatorname{Tr}}
\newcommand{\ip}[2]{\langle #1,#2\rangle}
\newcommand{\abs}[1]{\lvert #1\rvert}
\newcommand{\norm}[1]{\lVert #1\rVert}
\pagestyle{fancy}
\fancyhf{}
\fancyhead[L]{\small Sparse Grid Moments and Quadratic Closure}
\fancyhead[R]{\small Research draft}
\fancyfoot[C]{\thepage}
\renewcommand{\headrulewidth}{0.3pt}
\title{\textbf{Sparse Grid Moments and Quadratic Closure}\\[5pt]
\large Exact finite-field laws, sharp obstructions,\\
and a complete cube rank enumerator}
\author{Research draft prepared for the ProveIt project}
\date{7 October 2026}

\begin{document}
\maketitle
\thispagestyle{plain}
\begin{abstract}
A grid-moment lemma in the recent \texttt{openai/math} arithmetic-progression
manuscript asserts independence for sampling patterns with at most as many
axes as the degree of the chart polynomial. We study the exact algebraic
boundary of this mechanism in a self-contained finite-field model. A matrix
of low-order coordinate margins determines every relation among the sampled
polynomial values. For independent sums of degree-$d$ product blocks, the
joint law converges to the uniform measure on the orthogonal complement of
this relation code, with optimal exponential rate
$1-(1-Q^{-1})^{d-1}$ over a field of order $Q$. Normalized zero-level moments
therefore converge to $Q^{\dim K_d(S)}$, not automatically to one.

For quadratic blocks we prove a finer statement: the first relative
correction is $(Q-1)(|\cl_2(S)|-|S|)Q^{-n}$, where $\cl_2$ is closure in a
concrete represented matroid. We then determine the entire rank enumerator
of the seven-dimensional matrix space attached to the binary three-cube,
over every finite field, with separate odd- and even-characteristic
formulas. This gives exact finite-$n$ moments for both the seven-site
punctured cube and the full eight-site cube. For arbitrary two-axis
patterns, we prove that the exact convergence exponent is the graph girth
minus two; forests have exact independence, and all even cycles admit a
closed finite-dimensional formula. We also prove the sharp
$2^{d+1}$ minimum-support obstruction and a sharp variance lower bound for
arbitrary nonnegative cutoffs on the first obstructed cube. The arguments
use elementary algebra, finite Fourier analysis, and explicit counting;
none assumes the headline arithmetic-progression theorem in the source
repository.
\end{abstract}

\noindent\textbf{Status.} Written proofs and exact-arithmetic checks are supplied.
The manuscript was generated in an interactive ChatGPT research session and
has not been independently refereed or Lean-verified. The closure correction,
graph-girth law, and complete cube enumerator are the principal proposed original
contributions; their priority is not established. Classical ingredients are
identified explicitly.

\vspace{5pt}
\noindent\textbf{Keywords.} Relative counting; polynomial factors; finite fields;
grid margins; trades; matroid closure; rank enumerators; quadratic forms.

\newpage
\tableofcontents
\newpage

\section{Research target, source review, and scope}
\label{gm:sec:context}

\subsection{The intersection of the repositories}
The \texttt{openai/math} release inspected for this work contains a catalogue
of 722 manuscripts in 372 families, and explicitly distinguishes different
stages of verification. Its family 159 concerns claimed quasipolynomial
bounds for arithmetic progressions. The present work focuses on a single
local ingredient of that family, not on certification of its global
conclusion: the bounded-grid-moment argument in
\texttt{03-relative-detection.tex}, with source label
\texttt{ap:path:grid} \cite{oa,oa-grid}.

The relevant ProveIt research overview is
\emph{Local Quantitative Refinements of Gowers's Proof of Szemer\'edi's
Theorem}. It organizes local work on density transfer, polynomial phase
partitions, inverse steps, moment estimates, and progression or cube
counts. It also explicitly distinguishes a research manuscript's placement
from formal status \cite{proveit}. This makes a precise, independently
proved interface theorem more useful than an unexamined invocation of a
headline result. The broader classical setting is Gowers's quantitative
proof of Szemer\'edi's theorem \cite{gowers}.

In the selected OpenAI source, the polynomial layer has degree $j$, and the
sampling pattern has $q\leq j$ axes. The variables $x_{a,b}$ are independent,
and a site $\alpha$ is evaluated at $\sum_a x_{a,\alpha_a}$. The proof isolates
a Fourier coefficient by extracting a mixed monomial involving every
axis. The same section uses the resulting moments in a densification
argument. The relative-counting framework of Conlon, Fox, and Zhao requires
control of a finite collection of enlarged patterns and their subproducts,
not merely control of one-point densities \cite[\S\S2.2,6.2--6.3]{cfz}.

Our finite-field model exposes the exact relations that the coefficient
extraction is detecting. It also answers what happens when the number of
axes exceeds the degree, or when one removes a small number of sites from
an obstructed pattern.

\subsection{Precise questions and answers}
We address five local questions suggested by these sources.
\begin{enumerate}[label=(\roman*)]
\item Which relations among sites survive arbitrarily many independent
polynomial blocks? They are exactly the vectors with all coordinate
margins of order at most $d$ equal to zero
(Theorem~\ref{gm:thm:relations}).
\item What replaces the assertion that every normalized grid moment is
$1+o(1)$? The exact limiting constant is
$Q^{\dim K_d(S)}$, with an explicit, sharp exponential mixing rate for
product blocks (Theorems~\ref{gm:thm:mixing} and~\ref{gm:thm:moments}).
\item Does marginal normalization always remove the leading quadratic
error? No. Its remaining coefficient counts sites in the quadratic closure
that were not sampled (Theorem~\ref{gm:thm:closure}).
\item What determines the exact error rate within the two-axis regime?
It is the girth minus two, unless the pattern is a forest and the error
vanishes identically (Theorem~\ref{gm:thm:girth}).
\item Can the first quadratic obstruction be evaluated completely, at
finite dimension and in small characteristic? Yes: its full rank
enumerator is given in Theorem~\ref{gm:thm:cube-enumerator}.
\end{enumerate}
These are research questions formulated here, not claims that every one
was previously advertised as open in the literature.

\subsection{What is inherited, and what is proposed as new}
The principle that tensor relations govern polynomial configurations is
part of the true-complexity viewpoint of Gowers and Wolf \cite{gw}.
Coordinate-margin kernels are closely related to trades and low-degree
codes. In particular, the familiar minimum trade volume $2^d$ corresponds
to a support size $2^{d+1}$ when the positive and negative parts are
counted together; see Ghorbani, Kamali, Khosrovshahi, and Krotov
\cite{trades}. We prove the field-valued version needed here rather than
presenting this threshold as an unprecedented discovery. Fourier inversion,
convolution mixing, and quadratic character-sum evaluation are also
standard tools.

The main proposed contribution is the combination of an exact relation
criterion, sharp block-mixing estimates, and a geometric correction after
marginal normalization. The graph-girth theorem determines an exact convergence exponent for
every two-axis pattern, while the explicit odd- and even-characteristic cube
rank enumerators strengthen the asymptotic obstruction to a complete
finite-dimensional formula. Targeted literature searches did not locate
these specific closure, girth-law, and enumerator formulas; that limited search is
not a proof of novelty. The mathematical claims below stand on their
written proofs, independently of any priority claim.

\subsection{Limits of the comparison}
The source lemma explicitly assumes $q\leq j$. Our counterexamples beyond
that range do \emph{not} refute it. Nor does a finite-field proof establish
the source's torus-chart estimates, nilsequence complexity bounds, or
quasipolynomial progression theorem. We do not claim an improved global
bound for $r_k(N)$, a new proof of Szemer\'edi's theorem, or a completed Lean
formalization. The repository review was selective: the catalogue,
verification notices, the indicated grid argument, and the relevant
ProveIt overview were inspected, not all manuscripts or formal proofs.

\section{The coordinate-margin relation code}
\label{gm:sec:relations}

Let $\F$ be a finite field of order $Q$ and characteristic $p$. Fix a
nontrivial additive character $\psi:\F\to\mathbb C^\times$. All expectations
are uniform unless specified otherwise. Let $q\geq1$ be the number of axes,
let $A_1,\ldots,A_q$ be finite nonempty label sets, and write
\[
 \Omega=\prod_{a=1}^q A_a,\qquad m_a=|A_a|.
\]
A sampling pattern is a nonempty set $S\subseteq\Omega$, of size $s$.
For independent vectors $x_{a,b}$ in a fixed vector space over $\F$, set
\begin{equation}\label{gm:eq:site-form}
 L_\alpha(x)=\sum_{a=1}^q x_{a,\alpha_a}\qquad(\alpha\in S).
\end{equation}
Distinctness concerns the formal labels $\alpha$, not the random values of
$L_\alpha$. Actual collisions are included in every average.

\begin{definition}[Margin matrix and relation code]
For $d\geq0$, the matrix $M_d(S)$ has columns indexed by $S$. Its rows are
indexed by a set $I\subseteq[q]$, $|I|\leq d$, and a label tuple
$\beta\in\prod_{i\in I}A_i$. The entry in row $(I,\beta)$ and column $\alpha$
is $\1_{\alpha_I=\beta}$. The empty set contributes the all-one row. Define
\begin{equation}\label{gm:eq:code}
 K_d(S)=\ker_\F M_d(S),\quad r_d(S)=\rank_\F M_d(S),\quad
 H_d(S)=K_d(S)^\perp\subseteq\F^S.
\end{equation}
Thus $\lambda\in K_d(S)$ precisely when
\begin{equation}\label{gm:eq:margins}
 \sum_{\alpha\in S:\,\alpha_I=\beta}\lambda_\alpha=0
 \quad\text{for every }|I|\leq d\text{ and every }\beta.
\end{equation}
The orthogonal complement uses the ordinary $\F$-bilinear dot product.
\end{definition}

When no ambiguity is possible we write $K,r,H$ and $k=s-r=\dim K$.
The definition is finite linear algebra; it contains no asymptotic
hypothesis or analytic rank parameter.

\begin{definition}[$p$-reduced polynomial]
A polynomial in coordinates $z_1,\ldots,z_e$ is $p$-reduced if each
coordinate exponent in each monomial is less than $p$.
\end{definition}
Over a prime field this is the usual unique reduced representative of a
polynomial function. Over an extension field the condition is stronger
than having exponents less than $Q$; it is stated explicitly to exclude
Frobenius aliasing in the coefficient argument. Products of distinct
coordinates satisfy it over every finite field.

\begin{theorem}[Exact polynomial relations]\label{gm:thm:relations}
Let $P\in\F[z_1,\ldots,z_e]$ be $p$-reduced and of exact total degree
$d\geq1$. For $\lambda\in\F^S$, the following are equivalent:
\begin{enumerate}[label=(\alph*)]
\item $\lambda\in K_d(S)$;
\item $\sum_{\alpha\in S}\lambda_\alpha P(L_\alpha(x))$ is identically zero;
\item the degree-$d$ homogeneous component of this sum is zero.
\end{enumerate}
Identities may be read either as polynomial identities or as identities
of functions on the indicated finite vector spaces. If $\lambda\notin K_d(S)$,
the sum has exact degree $d$, rather than merely being nonconstant.
\end{theorem}

\begin{proof}
First expand a monomial of $P$ after substitution
\eqref{gm:eq:site-form}. A resulting monomial involves variables from at
most $d$ axes. If these axes form $I$, and the labels used on them are
$\beta$, its coefficient in the sum over $\alpha$ is its coefficient for a
single site multiplied by the margin in \eqref{gm:eq:margins}. At a given
axis there cannot be two different labels in the same site expansion.
Thus all coefficients vanish when $\lambda\in K_d(S)$, including the
constant coefficient because the empty margin vanishes.

For the converse, choose a nonzero top-degree monomial
$c z_1^{b_1}\cdots z_e^{b_e}$, where $\sum_j b_j=d$ and $0\leq b_j<p$.
Fix a nonempty $I\subseteq[q]$ with $|I|\leq\min(d,q)$, and prescribed labels
$\beta_i$ on $I$. Distribute the $d$ coordinate units among the axes of
$I$, using each axis at least once. Write the resulting nonnegative
integers as $b_{i,j}$, so that $\sum_{i\in I}b_{i,j}=b_j$.
The coefficient of
\[
 \prod_{i\in I}\prod_{j=1}^e x_{i,\beta_i,j}^{b_{i,j}}
\]
in the degree-$d$ part of the sum is
\begin{equation}\label{gm:eq:coefficient-extraction}
 c\prod_{j=1}^e\frac{b_j!}{\prod_{i\in I}b_{i,j}!}
 \sum_{\alpha\in S:\,\alpha_I=\beta}\lambda_\alpha.
\end{equation}
The multiplier is nonzero in $\F$: all factorial arguments are less than
$p$. No other original monomial contributes, since the total exponent in
base coordinate $j$ recovers $b_j$. Vanishing of the top component
therefore forces every nonempty margin of order at most $d$ to vanish.
Summing the singleton margins on any one axis gives the empty margin.

Finally, substitution does not increase the degree of any single
variable beyond the corresponding $b_j<p\leq Q$. A polynomial whose
individual degrees are less than $Q$ is zero as a function on $\F^N$ only
if all its coefficients vanish: apply the univariate root bound
successively to the variables. This proves the function version too.
\end{proof}

\begin{remark}
The argument uses products of coordinate factorials, not a division by
$d!$. It therefore remains valid when the total degree $d$ is at least the
characteristic, provided the stated coordinatewise hypothesis holds.
Without that hypothesis the margin degree need not be the displayed
formal degree. For example $P(z)=z^p$ is additive and has the relation
structure of degree one.
\end{remark}

\begin{proposition}[Full-grid rank]\label{gm:prop:full-rank}
For the full product $\Omega$,
\begin{equation}\label{gm:eq:full-rank}
 r_d(\Omega)=\sum_{\substack{I\subseteq[q]\\|I|\leq d}}
                      \prod_{i\in I}(m_i-1).
\end{equation}
In particular, $K_d(S)=0$ for every $S$ when $q\leq d$.
\end{proposition}

\begin{proof}
Choose a distinguished label in each $A_i$. The constant function and the
indicators of the other $m_i-1$ labels form a basis of the function space on
$A_i$. Tensor products of these bases form a basis of the function space
on $\Omega$. The row span of $M_d(\Omega)$ is exactly the span of those
tensors that use a nonconstant factor in at most $d$ coordinates. Counting
this subset of the basis proves \eqref{gm:eq:full-rank}. If $q\leq d$, the
rows with $I=[q]$ already isolate each site, including after restriction
to $S$.
\end{proof}

This proof does not divide by $m_i$, and hence remains valid when the
characteristic divides an axis size. An equivalent description of the
kernel uses the decomposition
\[
 \F^{A_i}=\F\delta_{b_i}\oplus
       \{u:\textstyle\sum_{a\in A_i}u(a)=0\}.
\]
Tensor products using a zero-sum factor on more than $d$ axes form a basis
of $K_d(\Omega)$. Indeed every margin of order at most $d$ sums out at
least one such factor, and the dimension is the complement of
\eqref{gm:eq:full-rank}.

\section{The first possible obstruction and its equality case}
\label{gm:sec:support}

\begin{theorem}[Sharp support threshold]\label{gm:thm:support}
Every nonzero vector in $K_d(\Omega)$ has at least $2^{d+1}$ nonzero
coordinates. Consequently,
\begin{equation}\label{gm:eq:small-independent}
 |S|<2^{d+1}\quad\Longrightarrow\quad K_d(S)=0.
\end{equation}
Whenever at least $d+1$ axes have two labels, the bound is attained.
\end{theorem}

\begin{proof}
We induct on $d$. For $d=0$, a nonzero vector with zero total sum has at
least two nonzero entries. For $d\geq1$, a nonzero kernel vector cannot be
supported at just one site. Choose an axis on which its support has at
least two different labels, and slice along that axis. Every nonzero
slice has all margins of order at most $d-1$ equal to zero on the remaining
axes, because these are margins of the original vector with the sliced
axis included. The induction hypothesis gives at least $2^d$ nonzero
entries in each such slice, and there are at least two slices.

Extend a vector on $S$ by zero to obtain the stated consequence. To attain
the bound, choose two labels on each of $d+1$ axes, hold all other axes
fixed, and assign the coefficient
\[
 \lambda_\varepsilon=(-1)^{\varepsilon_1+\cdots+\varepsilon_{d+1}}
 \qquad(\varepsilon\in\{0,1\}^{d+1}).
\]
Every margin involving at most $d$ axes sums over at least one of the
chosen pairs and vanishes. In characteristic two the two signs are equal,
but their sum is still zero.
\end{proof}

\begin{theorem}[Extremizers with exactly $d+1$ axes]
\label{gm:thm:extremizers}
Suppose $q=d+1$. Every nonzero vector in $K_d(\Omega)$ with support
$2^{d+1}$ is a nonzero scalar multiple of
\begin{equation}\label{gm:eq:extreme-tensor}
 (\delta_{a_1}-\delta_{b_1})\otimes\cdots\otimes
 (\delta_{a_q}-\delta_{b_q}),\qquad a_i\ne b_i.
\end{equation}
There are $\prod_i\binom{m_i}{2}$ such projective kernel vectors, or
$(Q-1)\prod_i\binom{m_i}{2}$ nonzero vectors.
\end{theorem}

\begin{proof}
In this case the kernel consists of arrays with every line sum zero, or
equivalently the tensor product of the $q$ zero-sum spaces. Slice along
the last axis. Every nonzero slice belongs to the corresponding kernel
on $q-1$ axes and has support at least $2^{q-1}$. Equality in the total
support bound forces exactly two slices. The last-axis line sums force
these slices to be negatives of one another, point by point. Induction
from $q=1$ proves \eqref{gm:eq:extreme-tensor}. Unordered pairs determine
the support, and all changes in the ordering of the pairs only change
the overall scalar.
\end{proof}

The support threshold is a trade-type bound, not an assertion of a new
minimum-distance phenomenon. Its use here is to give an exact,
characteristic-independent testing threshold for polynomial grid moments.
For more than $d+1$ axes, a complete classification of minimum-support
relations requires additional work; Theorem~\ref{gm:thm:extremizers} does
not assert that classification.

\section{Product blocks: exact limiting law and sharp mixing}
\label{gm:sec:mixing}

For $d,n\geq1$, define
\begin{equation}\label{gm:eq:product-block}
 B_{d,n}(z)=\sum_{\ell=1}^n\prod_{j=1}^d z_{j,\ell},
 \qquad z\in\F^{d n}.
\end{equation}
Apply the same polynomial at every site:
\[
 Y_n=\bigl(B_{d,n}(L_\alpha(x))\bigr)_{\alpha\in S}\in\F^S.
\]
The $n$ coordinate blocks are independent, so $Y_n$ is the sum of $n$
independent copies of $Y_1$. Fix $K=K_d(S)$, $H=H_d(S)$, and $r=\dim H$.

\begin{lemma}[Nonvanishing probability for multilinear maps]
\label{gm:lem:multilinear}
A nonzero $t$-linear map from a product of finite-dimensional $\F$-vector
spaces to another such space is nonzero on independent uniform inputs
with probability at least $(1-Q^{-1})^t$.
\end{lemma}

\begin{proof}
For $t=0$ the claim means that a fixed nonzero vector is nonzero. For the
induction step, fix vectors in the last $t-1$ inputs and a linear
functional on the output that witness a nonzero linear functional of the
first input. With probability $1-Q^{-1}$ this functional is nonzero, so
the contraction in the remaining $t-1$ inputs is a nonzero multilinear
map. Conditional on each such first input, apply the induction hypothesis.
\end{proof}

\begin{lemma}[One-block Fourier spectrum]\label{gm:lem:block-spectrum}
Put
\begin{equation}\label{gm:eq:theta}
 \theta=\theta_{Q,d}=1-(1-Q^{-1})^{d-1}.
\end{equation}
For $\lambda\in\F^S$, define
\[
 b_\lambda=\E\psi\!\left(\sum_{\alpha\in S}
                  \lambda_\alpha B_{d,1}(L_\alpha(x))\right).
\]
Then $b_\lambda=1$ on $K$, while $0\leq b_\lambda\leq\theta$ off $K$.
Moreover $b_{t e_\alpha}=\theta$ for every site $\alpha$ and every
$t\in\F^\times$.
\end{lemma}

\begin{proof}
The phase polynomial is $d$-linear in the $d$ groups of variables indexed
by the coordinate $j$ of the block. Theorem~\ref{gm:thm:relations} says
that it is zero exactly for $\lambda\in K$. Otherwise, integrating over
the last group gives the indicator that its vector of linear coefficients
is zero. Thus $b_\lambda$ is real and nonnegative. This coefficient vector
is a nonzero $(d-1)$-linear map; Lemma~\ref{gm:lem:multilinear} gives the
upper bound $\theta$. For $d=1$ it is a nonzero constant vector and the
bound is zero.

For $t e_\alpha$, the phase is $t$ times the product of $d$ independent
uniform scalars. The last scalar averages to zero unless one of the first
$d-1$ scalars is zero. This event has probability $\theta$.
\end{proof}

\begin{theorem}[Exact support space and optimal exponential rate]
\label{gm:thm:mixing}
The random vector $Y_n$ takes values in $H$, and for $n\geq1$,
\begin{equation}\label{gm:eq:tv}
 (1-Q^{-1})\theta^n
 \ \leq\ \norm{\Law(Y_n)-U_H}_{\TV}
 \ \leq\ \frac12\sqrt{Q^r-1}\,\theta^n,
\end{equation}
where $U_H$ denotes uniform measure on $H$. Also, for every $y\in H$,
\begin{equation}\label{gm:eq:pointwise}
 \left|Q^r\PP(Y_n=y)-1\right|\leq(Q^r-1)\theta^n.
\end{equation}
For $d\geq2$, the maximum absolute value of a nontrivial one-block Fourier
coefficient on $H$ is exactly $\theta$. For $d=1$, $Y_n$ is exactly uniform
on $H$ for every $n\geq1$.
\end{theorem}

\begin{proof}
Every relation in $K$ annihilates $Y_n$ by
Theorem~\ref{gm:thm:relations}. Characters of $H$ are indexed by
$\F^S/K$: two frequencies induce the same character precisely when their
difference lies in $K$. To see this also for extension fields, if a dot
product with an $\F$-linear space takes any nonzero value, scalar
multiplication makes it range over all of $\F$, on which $\psi$ is
nontrivial. Hence trace orthogonality and $\F$-orthogonality give the same
annihilator here.

Fourier inversion on $H$ gives
\begin{equation}\label{gm:eq:inversion}
 Q^r\PP(Y_n=y)=\sum_{\lambda\in\F^S/K}
                  b_\lambda^n\psi(-\ip\lambda y).
\end{equation}
The upper pointwise bound follows by the triangle inequality. Parseval
and Cauchy--Schwarz give the upper total-variation bound.

The coordinate projection $H\to\F$ at each site is surjective: otherwise
it is zero, forcing the singleton frequency $e_\alpha$ into $K$, contrary
to the empty-margin equation. Thus $U_H(Y_\alpha=0)=Q^{-1}$. The last
clause of Lemma~\ref{gm:lem:block-spectrum} gives the exact marginal
\begin{equation}\label{gm:eq:marginal}
 \PP(Y_{n,\alpha}=0)=Q^{-1}\bigl(1+(Q-1)\theta^n\bigr).
\end{equation}
The difference on this event is $(1-Q^{-1})\theta^n$, proving the lower
bound and optimality of the exponential rate. The same spectral argument
with $\theta=0$ proves the degree-one assertion.
\end{proof}

Define the zero-level weight and its exact mean by
\begin{equation}\label{gm:eq:weights}
 \nu_n(z)=Q\1_{B_{d,n}(z)=0},\qquad
 \mu_n=\E\nu_n=1+(Q-1)\theta^n,\qquad
 \bar\nu_n=\nu_n/\mu_n.
\end{equation}
Write $M_S(n)=\E\prod_{\alpha\in S}\nu_n(L_\alpha)$ and
$\bar M_S(n)=\E\prod_{\alpha\in S}\bar\nu_n(L_\alpha)$.

\begin{theorem}[Exact limiting moment and effective error]
\label{gm:thm:moments}
With $k=s-r$,
\begin{align}
 M_S(n)&=\sum_{\lambda\in\F^S}b_\lambda^n
       =Q^k\sum_{\lambda\in\F^S/K}b_\lambda^n,
       \label{gm:eq:raw-moment}\\
 Q^k&\leq M_S(n)\leq Q^k\bigl(1+(Q^r-1)\theta^n\bigr),
       \label{gm:eq:raw-bound}\\
 \left|\frac{\bar M_S(n)}{Q^k}-1\right|
     &\leq\bigl(Q^r-1+s(Q-1)\bigr)\theta^n.
       \label{gm:eq:normalized-bound}
\end{align}
In particular,
\begin{equation}\label{gm:eq:limiting-moment}
 \lim_{n\to\infty}\bar M_S(n)=Q^{\dim K_d(S)}.
\end{equation}
Thus the limit is one if and only if $K_d(S)=0$.
\end{theorem}

\begin{proof}
Expand each factor $Q\1_{u=0}=\sum_{t\in\F}\psi(tu)$ and use independence
of the $n$ blocks. All Fourier coefficients are nonnegative, so the
trivial quotient frequency supplies the lower bound. The upper bound is
Lemma~\ref{gm:lem:block-spectrum}. Finally, with $a=(Q-1)\theta^n$,
\[
 0\leq1-(1+a)^{-s}\leq sa.
\]
Divide \eqref{gm:eq:raw-bound} by $(1+a)^s$ and bound the two deviations
from one separately to obtain \eqref{gm:eq:normalized-bound}.
\end{proof}

\begin{corollary}[A finite family of tests]
\label{gm:cor:family}
Let $\mathcal S$ be a finite family of sampling patterns with
$K_d(S)=0$ for every $S\in\mathcal S$. For $d\geq2$, all their normalized
moments are within $\varepsilon$ of one whenever
\[
 n\geq
 \max_{S\in\mathcal S}
 \frac{\log\bigl((Q^{|S|}-1+|S|(Q-1))/\varepsilon\bigr)}{-\log\theta}.
\]
Integer rounding upward is understood. Every subpattern of an independent
pattern is independent, so this also controls all its subproducts.
\end{corollary}

\begin{proof}
Apply \eqref{gm:eq:normalized-bound} to each pattern. A dependence on a
subpattern would extend by zero to a dependence on the original pattern.
\end{proof}

\subsection{A genuinely sparse parameter regime}
At fixed $Q$, the level set has limiting density $Q^{-1}$; sparsity can be
made arbitrarily strong by allowing $Q$ to grow. The estimates remain
explicit in that regime. Since $r\geq1$ and
$\theta\leq(d-1)/Q$ for $d\geq2$,
\begin{equation}\label{gm:eq:large-field}
 \left|\frac{\bar M_S(n)}{Q^{s-r}}-1\right|
 \leq (s+1)(d-1)^n Q^{r-n}.
\end{equation}
For quadratics this becomes $(s+1)Q^{r-n}$. Consequently, for any fixed
independent pattern with $s$ sites, $n>s$ product blocks suffice for
asymptotic independence as $Q\to\infty$, while the density tends to zero.
This is a transparent model of a sparse testing problem. It is not a
claim that the constants in the torus-chart source become independent of
all chart parameters.

\section{Universality for other independent polynomial blocks}
\label{gm:sec:general-polynomial}

The product model gives sharp, positive Fourier coefficients. A weaker
but explicit universality statement holds for any $p$-reduced block.

\begin{lemma}[An elementary polynomial-bias bound]
\label{gm:lem:polynomial-bias}
Let $F$ be a $p$-reduced polynomial of exact degree $d\geq2$ over $\F$.
Then
\begin{equation}\label{gm:eq:xi}
 \left|\E_x\psi(F(x))\right|
 \leq \xi_{Q,d}:=
 \bigl(1-(1-Q^{-1})^{d-1}\bigr)^{1/2^{d-1}}<1.
\end{equation}
For a nonconstant affine polynomial the bias is zero.
\end{lemma}

\begin{proof}
The $d$-fold additive derivative of the top component is a nonzero
$d$-linear form $T$. For completeness, if a nonzero monomial is
$c z_1^{b_1}\cdots z_e^{b_e}$, take $b_j$ derivative vectors in the $j$th
coordinate direction. Its coefficient in $T$ is
$c\prod_j b_j!\ne0$. Distinct original monomials are separated by these
coordinate multiplicities.

After $d-1$ additive derivatives, the polynomial is affine in the base
variable $x$, with linear part $T(x,h_1,\ldots,h_{d-1})$. Repeated
Cauchy--Schwarz, equivalently the inequality between the mean and the
$(d-1)$st Gowers seminorm, gives
\[
 \left|\E_x\psi(F(x))\right|^{2^{d-1}}
 \leq \E_{h_1,\ldots,h_{d-1}}
       \left|\E_x\psi(\Delta_{h_1}\cdots\Delta_{h_{d-1}}F(x))\right|.
\]
The inner absolute value is the indicator that the linear part is zero.
Lemma~\ref{gm:lem:multilinear} bounds the probability of this event by
$1-(1-Q^{-1})^{d-1}$. Orthogonality proves the affine case.
\end{proof}

\begin{theorem}[Block universality]\label{gm:thm:universality}
Let $P$ be as in Theorem~\ref{gm:thm:relations}, and let
$P_n=\sum_{\ell=1}^n P(z^{(\ell)})$ on independent coordinate blocks.
The vector $(P_n(L_\alpha))_{\alpha\in S}$ converges to $U_{H_d(S)}$.
For $d\geq2$, its total-variation distance is at most
\[
 \frac12\sqrt{Q^{r_d(S)}-1}\,\xi_{Q,d}^n.
\]
For $d=1$ the distribution is exactly uniform after one block.
\end{theorem}

\begin{proof}
Theorem~\ref{gm:thm:relations} identifies the trivial characters and says
that every other phase has exact degree $d$. It is still $p$-reduced in
the site variables. Apply Lemma~\ref{gm:lem:polynomial-bias} and repeat the
Fourier--Parseval proof of Theorem~\ref{gm:thm:mixing}. A constant term in
$P$ does not alter the support space: every vector in $K_d(S)$ has zero
total coefficient.
\end{proof}

The exact relation space, rather than a special symmetry of product
blocks, therefore determines the limiting law throughout this class.
Only the sharp rate, positivity, and the later rank-enumerator formulas
are specific to the product model.

\section{Cutoff-independent obstructions on the first cube}
\label{gm:sec:cutoffs}

For $w:\F\to\mathbb C$, use the normalized transform
\[
 \widehat w(t)=Q^{-1}\sum_{u\in\F}w(u)\psi(-tu).
\]
For the product-block vector of Section~\ref{gm:sec:mixing}, Fourier
expansion immediately yields the following more general statement.

\begin{proposition}[Arbitrary test functions]\label{gm:prop:weighted}
For functions $w_\alpha:\F\to\mathbb C$,
\begin{align}
 \lim_{n\to\infty}\E\prod_{\alpha\in S}w_\alpha(Y_{n,\alpha})
 &=\E_{y\in H_d(S)}\prod_{\alpha\in S}w_\alpha(y_\alpha)\notag\\
 &=\sum_{\lambda\in K_d(S)}\prod_{\alpha\in S}\widehat w_\alpha(\lambda_\alpha).
 \label{gm:eq:weighted-limit}
\end{align}
The absolute difference from the limit is at most
\begin{equation}\label{gm:eq:weighted-error}
 \theta^n\sum_{\lambda\notin K_d(S)}
       \prod_{\alpha\in S}|\widehat w_\alpha(\lambda_\alpha)|
 \leq\theta^n\prod_{\alpha\in S}\norm{\widehat w_\alpha}_{\ell^1}.
\end{equation}
\end{proposition}

\begin{proof}
Insert Fourier inversion for every function. The expectation of the
resulting character is $b_\lambda^n$, which is one on the kernel and at
most $\theta^n$ in absolute value elsewhere. Uniform measure on $H_d(S)$
has exactly the same trivial characters.
\end{proof}

\begin{theorem}[A sharp variance obstruction]\label{gm:thm:variance}
Let $S=\{0,1\}^{d+1}$ and let $w:\F\to\mathbb R$ have mean one. Put
$V=\E(w-1)^2$ and $m=2^d$. Then
\begin{equation}\label{gm:eq:cube-cutoff}
 \lim_{n\to\infty}\E\prod_{\alpha\in S}w(Y_{n,\alpha})
   =1+\sum_{t\ne0}|\widehat w(t)|^{2m}
   \geq1+(Q-1)^{1-m}V^m.
\end{equation}
Equality in the inequality holds precisely when the nonzero Fourier
magnitudes are all equal. In the class of nonnegative mean-one $w$, the
lower bound is attained for every $0\leq V\leq Q-1$.
\end{theorem}

\begin{proof}
Proposition~\ref{gm:prop:full-rank} gives a one-dimensional relation
space, spanned by $\varepsilon_\alpha=(-1)^{|\alpha|}$. There are $m$ plus
and $m$ minus coefficients. Since $w$ is real,
$\widehat w(-t)=\overline{\widehat w(t)}$. Thus
\eqref{gm:eq:weighted-limit} is the asserted sum of $2m$th powers. In
characteristic two the coefficients coincide and the Fourier coefficients
of real $w$ are real, so the same formula holds.

Parseval gives $V=\sum_{t\ne0}|\widehat w(t)|^2$. The power-mean
inequality for $Q-1$ nonnegative numbers proves the bound and the equality
criterion. For $0\leq V\leq Q-1$, set
\begin{equation}\label{gm:eq:cutoff-extremizer}
 w(u)=1+a\bigl(Q\1_{u=0}-1\bigr),\qquad
 a=\sqrt{V/(Q-1)}.
\end{equation}
This is nonnegative and has every nonzero Fourier coefficient equal to
$a$. It has variance $V$ and attains the lower bound. Conversely,
nonnegative mean-one functions satisfy $\E w^2\leq Q$, since their
nonnegative values sum to $Q$, so $V\leq Q-1$ is the full possible range.
\end{proof}

Thus a nonconstant real cutoff cannot restore independence on this first
obstructed cube by increasing the number of blocks. Its variance forces
a positive excess. For the zero-level cutoff $w=Q\1_0$, the limit is $Q$.
Marginal normalization of $w(B_{d,n})$ does not change this limiting value,
since its marginal mean tends to $\E w=1$.

\section{Quadratic closure and the first normalized correction}
\label{gm:sec:closure}

Now set $d=2$. Put $N=\sum_a m_a$, and represent each site by its one-hot
vector $v_\alpha\in\F^N$: in the block indexed by axis $a$, it has a one
at label $\alpha_a$ and zeros elsewhere. Define
\begin{equation}\label{gm:eq:matrix-space}
 A_\lambda=\sum_{\alpha\in S}\lambda_\alpha v_\alpha v_\alpha^{\mathsf T},
 \qquad
 \mathcal V_S=\Span_\F\{v_\alpha v_\alpha^{\mathsf T}:\alpha\in S\}.
\end{equation}
These are symmetric matrices, but the character sum is bilinear in two
\emph{independent} groups of variables. No division by two occurs.

\begin{lemma}[The matrix representation]\label{gm:lem:matrix-representation}
The kernel of $\lambda\mapsto A_\lambda$ is $K_2(S)$, so
$\dim\mathcal V_S=r_2(S)$. For one quadratic block,
\begin{equation}\label{gm:eq:rank-bias}
 b_\lambda=Q^{-\rank A_\lambda}.
\end{equation}
\end{lemma}

\begin{proof}
Diagonal entries of $A_\lambda$ are singleton margins, off-diagonal
entries from distinct axes are pair margins, and off-diagonal entries
within one axis are zero. The empty margin is the sum of the singleton
margins on any axis. This proves the kernel assertion. The phase of a
single block is $x^{\mathsf T}A_\lambda y$ with independent uniform
$x,y\in\F^N$. Averaging over $x$ gives $\1_{A_\lambda y=0}$, whose mean
is $Q^{-\rank A_\lambda}$.
\end{proof}

\begin{definition}[Quadratic closure and rank enumerator]
With the ambient full grid $\Omega$ fixed, put
\begin{align}
 \cl_2(S)&=\{\beta\in\Omega:v_\beta v_\beta^{\mathsf T}\in\mathcal V_S\},
                 \label{gm:eq:closure-definition}\\
 R_S(z)&=\sum_{A\in\mathcal V_S}z^{\rank A}.
                 \label{gm:eq:rank-enumerator}
\end{align}
The first is the ordinary closure operation of the matroid represented by
the site matrices. Write $c=|\cl_2(S)|$.
\end{definition}

\begin{theorem}[Rank-one rigidity]\label{gm:thm:rank-one}
Every nonzero rank-one matrix in $\mathcal V_\Omega$ has the form
$t v_\beta v_\beta^{\mathsf T}$ for a unique $\beta\in\Omega$ and unique
$t\in\F^\times$. Consequently $\mathcal V_S$ has exactly $(Q-1)c$
rank-one matrices.
\end{theorem}

\begin{proof}
A nonzero symmetric rank-one matrix over any field is $tww^{\mathsf T}$
with $t\ne0$ and $w\ne0$. Indeed writing it as an outer product and using
symmetry makes its two factors proportional. Inside any axis block, the
off-diagonal entries vanish. Hence $w$ has at most one nonzero coordinate
in each block.

Let $u_a$ be the indicator vector of the entire axis-$a$ block. Every
site matrix satisfies
$(v_\alpha v_\alpha^{\mathsf T})u_a=v_\alpha$, independently of $a$.
Therefore $A u_a$ is independent of $a$ for every $A\in\mathcal V_\Omega$.
For $A=tww^{\mathsf T}$, this says that the coordinate sum of $w$ in each
block is the same scalar $h$. If $h=0$, the at-most-one-nonzero-coordinate
property forces all blocks to be zero, a contradiction. Thus each block
has exactly one nonzero coordinate, and its value is $h$. It follows that
$w=h v_\beta$ for a unique site $\beta$. Absorbing $h^2$ into $t$ proves
the assertion. The converse is the definition of closure.
\end{proof}

\begin{theorem}[Exact moment and closure correction]\label{gm:thm:closure}
Let $s=|S|$, $r=r_2(S)$, $k=s-r$, $c=|\cl_2(S)|$, and $z=Q^{-n}$.
Then for every $n\geq1$,
\begin{equation}\label{gm:eq:rank-moment}
 M_S(n)=Q^k R_S(z),\qquad
 \bar M_S(n)=Q^k\frac{R_S(z)}{(1+(Q-1)z)^s}.
\end{equation}
As $n\to\infty$ with $Q,S$ fixed,
\begin{equation}\label{gm:eq:closure-correction}
 \frac{\bar M_S(n)}{Q^k}
   =1+(Q-1)(c-s)Q^{-n}+O_{Q,S}(Q^{-2n}).
\end{equation}
An explicit remainder bound, valid for every $0\leq z\leq1$, is
\begin{equation}\label{gm:eq:closure-remainder}
 \left|\frac{R_S(z)}{(1+az)^s}-1-a(c-s)z\right|
 \leq C_{Q,S}z^2,\qquad a=Q-1,
\end{equation}
where one may take
\[
 C_{Q,S}=Q^r-1-ac+a^2cs+	frac12(1+ac)s(s+1)a^2.
\]
\end{theorem}

\begin{proof}
Every matrix in $\mathcal V_S$ has $Q^k$ frequency preimages. Substitute
\eqref{gm:eq:rank-bias} in \eqref{gm:eq:raw-moment}; the marginal mean is
$1+(Q-1)Q^{-n}$. This proves the exact formula.
Theorem~\ref{gm:thm:rank-one} gives
\[
 R_S(z)=1+ac z+T(z),\qquad
 0\leq T(z)\leq(Q^r-1-ac)z^2\quad(0\leq z\leq1).
\]
Taylor's theorem applied to $(1+az)^{-s}$ gives
\[
 (1+az)^{-s}=1-sa z+E(z),\qquad
 0\leq E(z)\leq\tfrac12s(s+1)a^2z^2.
\]
Multiply the two expressions, use $(1+az)^{-s}\leq1$, and bound
$1+ac z\leq1+ac$. The cross term has magnitude $a^2csz^2$.
This gives \eqref{gm:eq:closure-remainder}, and the asymptotic statement
follows.
\end{proof}

The coefficient is nonnegative because $S\subseteq\cl_2(S)$. It is not a
function of $r$ and $s$ alone: the locations of unsampled rank-one points
in the span matter. In particular, the first correction vanishes exactly
when $S$ is quadratically closed in the chosen ambient grid.

\begin{corollary}[Six sites versus seven]\label{gm:cor:six-seven}
Every pattern with at most six sites is quadratically independent and
closed. Its normalized moment is $1+O_{Q,S}(Q^{-2n})$.
A seven-site pattern can be independent but not closed: deleting one
vertex from the binary three-cube gives $r=s=7$ and $c=8$, hence
\[
 \bar M_S(n)=1+(Q-1)Q^{-n}+O_Q(Q^{-2n}).
\]
For the full eight-site cube, $r=7$, $k=1$, and $c=s=8$, so
$\bar M_S(n)=Q+O_Q(Q^{1-2n})$.
\end{corollary}

\begin{proof}
If a hidden site belongs to the closure of a pattern of size at most six,
there is a dependence on at most seven sites, contradicting
Theorem~\ref{gm:thm:support} with $d=2$. The same theorem gives independence.
The full binary three-cube has rank $1+3+3=7$ and one relation, all of whose
eight coefficients are nonzero. Deleting any vertex leaves seven
independent columns spanning the same space. Theorem~\ref{gm:thm:closure}
now applies.
\end{proof}

\section{Two-axis patterns: the exact rate is determined by girth}
\label{gm:sec:girth}

For $q=2$, regard $S\subseteq A_1\times A_2$ as the edge set of a simple
bipartite graph $G$. Every such pattern is quadratically independent and
closed, because a pair margin already isolates an edge. Nevertheless,
the normalized error need not have order $Q^{-2n}$. The graph determines
its exact order.

Let $s=|E(G)|$, $a=Q-1$, and $z=Q^{-n}$. Write $g(G)$ for the length of a
shortest cycle when $G$ contains a cycle. A bipartite graph has even girth
$g\geq4$. Vertices of degree zero may be discarded.

\begin{theorem}[Girth law for quadratic grid moments]
\label{gm:thm:girth}
If $G$ is a forest, then $\bar M_G(n)=1$ for every $n\geq1$.
If $G$ contains a cycle and has girth $g$, then $\bar M_G(n)>1$ for every
$n\geq1$, and there is a positive integer $C_{G,Q}$ such that
\begin{equation}\label{gm:eq:girth-asymptotic}
 \bar M_G(n)=1+C_{G,Q}Q^{-(g-2)n}
                +O_{G,Q}(Q^{-(g-1)n}).
\end{equation}
Let $c_g(G)$ be the number of unoriented simple cycles of length $g$, and
put
\begin{equation}\label{gm:eq:cycle-zero-count}
 N_{g,Q}=\frac{(Q-1)^g+(Q-1)}{Q}.
\end{equation}
Then $C_{G,Q}\geq c_g(G)N_{g,Q}>0$. More explicitly, for every $n\geq1$,
\begin{equation}\label{gm:eq:girth-finite-bounds}
 \frac{c_g(G)N_{g,Q}z^{g-2}(1-z^2)}{(1+az)^s}
 \leq\bar M_G(n)-1\leq Q^s z^{g-2}.
\end{equation}
Thus the exponent $g-2$ is exact, rather than merely a permissible upper
bound for the error.
\end{theorem}

\begin{proof}
Change the signs on the coordinates in the second bipartition class.
This congruence sends a site matrix to
$(e_u-e_v)(e_u-e_v)^{\mathsf T}$. Hence the matrix $A_\lambda$ becomes the
weighted graph Laplacian
\[
 \mathcal L_\lambda=D\diag(\lambda)D^{\mathsf T},
\]
where $D$ is an oriented incidence matrix. The orientation of an edge
does not change its outer product. For a frequency $\lambda$, let
$F=\supp\lambda$ and $t=|F|$.

If $F$ is a forest, its incidence columns are independent, as follows by
removing leaves. Since all retained weights are nonzero,
$\rank\mathcal L_\lambda=t$. If $F$ contains a cycle, its incidence rank
is $|V(F)|-c(F)<t$, where $c(F)$ counts nonempty connected components.
Thus $\rank\mathcal L_\lambda<t$ in this case. Consequently
\begin{equation}\label{gm:eq:girth-difference}
 R_G(z)-(1+az)^s
  =\sum_{\lambda:\,\supp\lambda\text{ contains a cycle}}
      \left(z^{\rank\mathcal L_\lambda}-z^{|\supp\lambda|}\right).
\end{equation}
Each summand is positive for $0<z<1$. This already proves the forest
assertion and strict positivity for cyclic graphs.

For a lower bound on the rank of a cyclic frequency, take a shortest
cycle in its support, of length $h\geq g$. This cycle is chordless within
the support. The corresponding $h\times h$ principal submatrix of
$\mathcal L_\lambda$ has nonzero entries along the cycle edges and arbitrary
diagonal entries. A vector in its kernel is determined by its first two
coordinates: the successive internal row equations determine the third
coordinate, then the fourth, and so on. Thus its nullity is at most two,
and
\begin{equation}\label{gm:eq:girth-rank-bound}
 \rank\mathcal L_\lambda\geq h-2\geq g-2.
\end{equation}
The upper finite bound now follows from \eqref{gm:eq:girth-difference},
with at most $Q^s$ frequencies and denominator $(1+az)^s\geq1$.

To attain rank $g-2$, support the weights on a fixed shortest cycle.
Orient its edges cyclically. The equation $\mathcal L_\lambda y=0$ says
that the currents $\lambda_i(y_i-y_{i+1})$ are all equal. Thus its kernel
has dimension two exactly when
\begin{equation}\label{gm:eq:cycle-reciprocal}
 \sum_{i=1}^{g}\lambda_i^{-1}=0;
\end{equation}
otherwise its kernel has dimension one. All weights are required to be
nonzero. Inversion permutes $\F^\times$, and additive-character
orthogonality counts the solutions of \eqref{gm:eq:cycle-reciprocal} as
\[
 Q^{-1}\left((Q-1)^g+(Q-1)(-1)^g\right)=N_{g,Q}.
\]
Here $g$ is even, so this number is positive over every finite field.
Different shortest cycles give different frequency supports. Their
contribution to \eqref{gm:eq:girth-difference} is at least
$c_g(G)N_{g,Q}(z^{g-2}-z^g)$, proving the lower finite bound.

Finally, no negative term on the right of
\eqref{gm:eq:girth-difference} has degree below $g$. Its coefficient of
$z^{g-2}$ is precisely the number of cyclic-support frequencies of rank
$g-2$; call this positive integer $C_{G,Q}$. Dividing by $(1+az)^s$, whose
constant term is one, gives \eqref{gm:eq:girth-asymptotic}.
\end{proof}

The leading coefficient need not count only shortest-cycle-supported
frequencies; larger supports may also yield rank $g-2$. The theorem makes
no assertion excluding them. Its exact coefficient description is
\[
 C_{G,Q}=\#\{\lambda\in\F^{E(G)}:\supp\lambda\text{ is cyclic and }
                              \rank\mathcal L_\lambda=g-2\}.
\]

\begin{corollary}[Every even cycle, exactly]\label{gm:cor:cycle-exact}
For the cycle $C_\ell$ of even length $\ell\geq4$, let
$N=N_{\ell,Q}=((Q-1)^\ell+(Q-1))/Q$. Then
\begin{align}
 R_{C_\ell}(z)
   &=(1+az)^\ell+N z^{\ell-2}+(a^\ell-N)z^{\ell-1}-a^\ell z^\ell,
          \label{gm:eq:cycle-enumerator}\\
 \bar M_{C_\ell}(n)
   &=1+\frac{N z^{\ell-2}+(a^\ell-N)z^{\ell-1}-a^\ell z^\ell}
                  {(1+az)^\ell}.
          \label{gm:eq:cycle-moment}
\end{align}
In particular, $C_{C_\ell,Q}=N$. The exact convergence exponent can be
arbitrarily large as the cycle length grows.
\end{corollary}

\begin{proof}
Every proper edge subset is a forest and contributes the baseline rank
equal to its size. Among the $a^\ell$ full-support frequencies, exactly
$N$ have rank $\ell-2$, by \eqref{gm:eq:cycle-reciprocal}; the rest have
rank $\ell-1$. Replace the baseline full-support term
$a^\ell z^\ell$ accordingly and normalize.
\end{proof}

For a rectangular four-cycle this gives
\[
 R_{C_4}(z)=1+4az+a(Q^2+3Q-3)z^2+Qa(Q^2-2)z^3,
\]
and the first normalized correction is
$a(Q^2-3Q+3)Q^{-2n}$. For a six-cycle it starts at $Q^{-4n}$, not
$Q^{-2n}$. Over $\F_2$, the exact cycle formula simplifies to
\[
 \bar M_{C_\ell}(n)
    =1+\frac{z^{\ell-2}(1-z)}{(1+z)^{\ell-1}}.
\]
For example, at $n=1$ the four-cycle moment is $28/27$ and the six-cycle
moment is $244/243$. The two-axis source regime therefore contains a
whole hierarchy of exact convergence rates, even though every pattern
has the same independent limiting law.

\section{A complete cube rank enumerator over every finite field}
\label{gm:sec:enumerator}

Let $\Omega=\{0,1\}^3$. The full cube and each one-vertex-deleted cube
span the same seven-dimensional matrix space. It is convenient to replace
the six-coordinate one-hot vector by
\[
 w_\alpha=(1,\alpha_1,\alpha_2,\alpha_3)^{\mathsf T}.
\]
There is a $6\times4$ matrix of full column rank carrying $w_\alpha$ to
$v_\alpha$. Therefore the replacement preserves ranks of the corresponding
matrix sums. Their space is
\begin{equation}\label{gm:eq:cube-space}
 \mathcal C_Q=
 \left\{
 \begin{pmatrix}
 a&b&c&d\\ b&b&e&f\\ c&e&c&g\\ d&f&g&d
 \end{pmatrix}: a,b,c,d,e,f,g\in\F
 \right\}.
\end{equation}
Indeed the seven functions $1,\alpha_i,\alpha_i\alpha_j$ are linearly
independent on the Boolean cube over every field. Thus all seven
parameters occur independently.

\begin{theorem}[Full rank distribution]\label{gm:thm:cube-enumerator}
Write
\[
 R_Q(z)=\sum_{A\in\mathcal C_Q}z^{\rank A}
        =1+a_1z+a_2z^2+a_3z^3+a_4z^4.
\]
In every characteristic, $a_1=8(Q-1)$. In odd characteristic,
\begin{align}
 a_2&=(Q-1)(Q^3+7Q^2+Q-7),\notag\\
 a_3&=Q^2(Q-1)(Q^3+Q^2+7Q-13),\label{gm:eq:odd-enumerator}\\
 a_4&=Q^2(Q-1)(Q^4-7Q+7).\notag
\end{align}
In characteristic two,
\begin{align}
 a_2&=(Q-1)(Q^3+8Q^2+Q-7),\notag\\
 a_3&=Q^2(Q-1)(Q^3+Q^2+8Q-14),\label{gm:eq:even-enumerator}\\
 a_4&=Q^2(Q-1)^2(Q^3+Q^2+Q-7).\notag
\end{align}
These formulas hold for every finite field, not just prime fields.
\end{theorem}

We give a counting proof rather than interpolation from numerical data.
The rank-one count follows from Theorem~\ref{gm:thm:rank-one}. Two
additional counts determine the remaining coefficients: a rank-two Gram
count and a first kernel moment.

\subsection{Why Gram counting counts each matrix exactly once}
A rank-two symmetric $4\times4$ matrix can be represented as
\[
 A_{ij}=B(u_i,u_j),\qquad u_0,u_1,u_2,u_3\in\F^2,
\]
where $B$ is a nondegenerate symmetric bilinear form and the four vectors
span $\F^2$. To construct this representation, quotient $\F^4$ by the
radical of its matrix form. Two spanning tuples give the same matrix
exactly when an isometry carries one tuple to the other: the map on their
spans is well defined because the kernels of the two Gram maps agree.
The isometry action on spanning tuples is free. Consequently one counts
spanning tuples for each isometry type and divides by its isometry-group
order.

\subsection{Rank two in odd characteristic}
When $Q$ is odd, replace $w_\alpha$ by
$(1,2\alpha_1-1,2\alpha_2-1,2\alpha_3-1)$. This is an invertible change of
coordinates. The transformed space consists exactly of symmetric
$4\times4$ matrices with all four diagonal entries equal. Hence the Gram
vectors must all have a common norm $t$.

A nondegenerate plane has two isometry types. The split type may be written
with norm $xy$; it has $2Q-1$ vectors of norm zero and $Q-1$ of each nonzero
norm. Its isometry group has order $2(Q-1)$. The anisotropic type is a
nonzero scalar multiple of the norm of $\F_{Q^2}/\F_Q$; it has one vector
of norm zero and $Q+1$ of each nonzero norm, and its isometry group has
order $2(Q+1)$. These assertions follow directly as follows. In the split
case the zero set is the union of the two isotropic lines, and an isometry
scales these lines reciprocally or interchanges them. In the anisotropic
case the norm map on the cyclic multiplicative group has kernel of order
$Q+1$; after choosing the image of a unit vector, an isometry fixing it is
identity or conjugation. These two cases exhaust nondegenerate planes by
diagonalization and whether a nonzero isotropic vector exists.

For nonzero $t$, if the norm fiber has $N_t$ vectors, each line meets it
in zero or two points. The number of ordered quadruples in that fiber
which fail to span the plane is therefore $8N_t$. In the split zero
fiber, the nonspanning count is $2Q^4-1$, the union of the two isotropic
lines counted with their shared zero tuple corrected. In the anisotropic
zero fiber there are no spanning tuples. It follows that
\begin{align*}
 a_2={}&\frac{(2Q-1)^4-2Q^4+1+(Q-1)((Q-1)^4-8(Q-1))}{2(Q-1)}\\
 &+\frac{(Q-1)((Q+1)^4-8(Q+1))}{2(Q+1)}\\
 ={}&(Q-1)(Q^3+7Q^2+Q-7).
\end{align*}
This proves the odd-characteristic rank-two formula.

\subsection{Rank two in characteristic two}
Use the original space \eqref{gm:eq:cube-space}. The Gram constraints are
\begin{equation}\label{gm:eq:even-gram-constraint}
 B(u_i,u_i)=B(u_0,u_i),\qquad 1\leq i\leq3.
\end{equation}
A nondegenerate symmetric bilinear plane in characteristic two has two
isometry types: alternating, and nonalternating. For the latter, an
orthogonal basis and the existence of square roots in a finite field
identify the form with the identity matrix. Their isometry-group orders
are respectively $Q(Q^2-1)$ and $Q$. The first is
$|\operatorname{SL}_2(\F)|$; for the second, a unit column $(a,b)$ satisfies
$a+b=1$ and determines its orthogonal unit companion $(b,a)$, giving $Q$
choices.

For the alternating plane, when $u_0\ne0$, condition
\eqref{gm:eq:even-gram-constraint} forces every $u_i$ into the line
$u_0^\perp=\F u_0$, so no tuple spans. When $u_0=0$, the last three vectors
are arbitrary. The nonspanning triples are the zero triple and the
nonzero triples contained in one of the $Q+1$ lines. The contribution is
\begin{equation}\label{gm:eq:alternating-count}
 \frac{Q^6-1-(Q+1)(Q^3-1)}{Q(Q^2-1)}=Q^3-1.
\end{equation}

For the nonalternating plane, choose coordinates $(u,t)$ with
\[
 B((U,T),(u,t))=Ut+Tu+Tt,\qquad B((u,t),(u,t))=t^2.
\]
Fix $u_0=(U,T)$. If $T\ne0$, the equation
$t^2=Tu+(U+T)t$ has exactly $Q$ solutions. Its intersection with the line
$\F u_0$ consists of $0$ and $u_0$. Thus there are $Q^3-8$ choices of
$(u_1,u_2,u_3)$ which span with $u_0$, and $Q(Q-1)$ such $u_0$.
If $T=0$ and $U\ne0$, the solution set consists of the two parallel lines
$t=0$ and $t=U$. Its intersection with $\F u_0$ is one entire line, giving
$(2Q)^3-Q^3=7Q^3$ spanning triples. There are $Q-1$ such $u_0$.
For $u_0=0$, all allowed vectors lie on the isotropic line, so there is
no contribution. Dividing by the isometry-group order $Q$, and adding
\eqref{gm:eq:alternating-count}, gives
\begin{align*}
 a_2&=(Q^3-1)
    +\frac{Q(Q-1)(Q^3-8)+(Q-1)7Q^3}{Q}\\
    &=(Q-1)(Q^3+8Q^2+Q-7).
\end{align*}

\subsection{A kernel moment determines ranks three and four}
For any seven-dimensional matrix space $\mathcal C$,
\begin{equation}\label{gm:eq:kernel-double-count}
 \sum_{A\in\mathcal C}|\ker A|
     =\sum_{x\in\F^4}Q^{7-\rho(x)},
 \qquad \rho(x)=\rank(A\mapsto Ax).
\end{equation}
We compute $\rho(x)$ by finding the annihilator of the image.

In odd characteristic use the equal-diagonal model. A vector $y$ is in
this annihilator exactly when
\[
 \sum_i x_i y_i=0,\qquad x_i y_j+x_j y_i=0\quad(i<j).
\]
If $x$ has at least three nonzero coordinates, the ratios $y_i/x_i$
on its support have pairwise sums zero, forcing them all to vanish;
coordinates off the support vanish too. A one-coordinate support also
has zero annihilator. For a two-coordinate support, the annihilator is
one-dimensional exactly when the two nonzero coordinates have equal
squares. There are $\binom42\cdot2(Q-1)=12(Q-1)$ such nonzero $x$.
Thus $\rho(x)=3$ on these vectors and is $4$ on all other nonzero vectors.

In characteristic two use \eqref{gm:eq:cube-space}. Write
$x=(x_0,x_1,x_2,x_3)$ and similarly for $y$. The annihilator equations are
\begin{align*}
 x_0y_0&=0,\\
 x_i y_0+(x_0+x_i)y_i&=0\quad(1\leq i\leq3),\\
 x_i y_j+x_j y_i&=0\quad(1\leq i<j\leq3).
\end{align*}
If $x_0\ne0$, then $y_0=0$. A nonzero annihilator exists exactly when the
tail is a nonempty vector with entries in $\{0,x_0\}$; it is then
one-dimensional. If $x_0=0$, the last equations make the tail of $y$
proportional to the nonzero tail of $x$. The remaining equations permit
a nonzero proportionality factor exactly when all nonzero tail entries
of $x$ are equal. Again the annihilator is one-dimensional. There are
$7(Q-1)$ vectors of each kind, for a total of $14(Q-1)$. Every other
nonzero vector has zero annihilator.

Set $\eta=12$ in odd characteristic and $\eta=14$ in characteristic two.
Equation~\eqref{gm:eq:kernel-double-count} becomes
\begin{equation}\label{gm:eq:kernel-moment-value}
 T_Q=Q^7+\eta(Q-1)Q^4+\bigl(Q^4-1-\eta(Q-1)\bigr)Q^3.
\end{equation}
On the other hand,
\[
 T_Q=Q^4+a_1Q^3+a_2Q^2+a_3Q+a_4,
 \qquad Q^7=1+a_1+a_2+a_3+a_4.
\]
Consequently
\begin{equation}\label{gm:eq:solve-rank-three}
 (Q-1)a_3=T_Q-Q^7-(Q^4-1)-(Q^3-1)a_1-(Q^2-1)a_2.
\end{equation}
Substitution of the two rank-two formulas and then use of the total count
proves \eqref{gm:eq:odd-enumerator} and \eqref{gm:eq:even-enumerator}.
This completes the proof of Theorem~\ref{gm:thm:cube-enumerator}.
\qed

\section{Exact seven- and eight-site moments and consequences}
\label{gm:sec:consequences}

Let $S_8=\{0,1\}^3$ and $S_7=S_8\setminus\{\alpha_0\}$ for any vertex
$\alpha_0$. In the quadratic product-block model, put $z=Q^{-n}$.

\begin{corollary}[Finite-dimensional formulas]\label{gm:cor:exact-seven-eight}
For every finite field and every $n\geq1$,
\begin{equation}\label{gm:eq:seven-eight}
 \bar M_{S_7}(n)=\frac{R_Q(z)}{(1+(Q-1)z)^7},\qquad
 \bar M_{S_8}(n)=\frac{Q R_Q(z)}{(1+(Q-1)z)^8},
\end{equation}
where $R_Q$ is the explicit polynomial of
Theorem~\ref{gm:thm:cube-enumerator}. In particular,
\begin{equation}\label{gm:eq:seven-eight-ratio}
 \bar M_{S_8}(n)=\frac{Q\bar M_{S_7}(n)}{1+(Q-1)Q^{-n}}.
\end{equation}
\end{corollary}

\begin{proof}
Both patterns span $\mathcal C_Q$. Their relation dimensions are zero and
one respectively. Apply \eqref{gm:eq:rank-moment}.
\end{proof}

Define
\begin{equation}\label{gm:eq:second-coefficient}
 b_Q=a_2-28(Q-1)^2
 =\begin{cases}
 (Q-1)(Q^3+7Q^2-27Q+21),&Q\text{ odd},\\
 (Q-1)(Q^3+8Q^2-27Q+21),&Q\text{ even}.
 \end{cases}
\end{equation}
Expanding the exact rational expressions gives
\begin{align}
 \bar M_{S_7}(n)&=1+(Q-1)Q^{-n}+b_QQ^{-2n}+O_Q(Q^{-3n}),
                  \label{gm:eq:seven-expansion}\\
 Q^{-1}\bar M_{S_8}(n)&=1+b_QQ^{-2n}+O_Q(Q^{-3n}).
                  \label{gm:eq:eight-expansion}
\end{align}
The coefficient $b_Q$ is positive for every prime power $Q\geq2$, as is
immediate from the displayed polynomials (the odd case starts at $Q=3$).
Thus the error orders shown here are genuinely attained, not merely
upper bounds.

\begin{table}[ht]
\centering
\caption{Exact cube rank counts; the last column is the coefficient in
\eqref{gm:eq:second-coefficient}. Each row was also checked by exhaustive
finite-field matrix enumeration.}
\label{gm:tab:ranks}
\begin{tabular}{rrrrrrr}
\toprule
$Q$&$a_0$&$a_1$&$a_2$&$a_3$&$a_4$&$b_Q$\\
\midrule
2&1&8&35&56&28&7\\
3&1&16&172&792&1,206&60\\
4&1&24&567&4,704&11,088&315\\
5&1&32&1,192&17,200&59,700&744\\
7&1&48&4,116&125,832&693,546&3,108\\
\bottomrule
\end{tabular}
\end{table}

For example,
\[
 R_2(z)=1+8z+35z^2+56z^3+28z^4.
\]
At $n=1$ one obtains $\bar M_{S_7}=320/243$ and
$\bar M_{S_8}=1280/729$, as checked independently by direct counting.
The eight-site value is below its limiting value $2$, whereas
\eqref{gm:eq:eight-expansion} places it above $2$ for all sufficiently
large $n$. Thus monotonicity in $n$ does not hold in this example.
There is no contradiction: the expansion describes the neighborhood of
$z=0$, not the whole interval $0<z\leq1/2$.

\subsection{Larger grids}
For a binary $q$-axis grid,
\begin{equation}\label{gm:eq:binary-limit}
 \lim_{n\to\infty}\bar M_{\{0,1\}^q}(n)
   =Q^{\,2^q-\sum_{i=0}^{\min(d,q)}\binom qi}.
\end{equation}
For quadratic blocks, the four-axis and five-axis limits are respectively
$Q^5$ and $Q^{16}$. These excesses are algebraic and persist even as the
nontrivial block biases tend to zero. A condition that asks all of these
moments to approach one is therefore incompatible with this model.

Conversely, the number of axes alone is not decisive for a partial grid:
Theorem~\ref{gm:thm:support} guarantees independence below $2^{d+1}$ sites,
however many axes are used. The exact matrix $M_d(S)$ decides all remaining
cases. This refines a sufficient condition based only on $q\leq d$ into
an if-and-only-if criterion for the present model.

\subsection{What this contributes to a densification audit}
For any finite collection of distinct-site patterns generated by a
Cauchy--Schwarz or densification argument, one can first compute their
margin ranks. An independent pattern passes the product-block test, with
the explicit error in \eqref{gm:eq:normalized-bound}. A dependent pattern
has a known limiting excess, and Theorem~\ref{gm:thm:variance} can rule out
repair by merely choosing a smoother nonconstant cutoff on the first
cube. At quadratic degree, closure distinguishes whether marginal
normalization improves the error from $Q^{-n}$ to $Q^{-2n}$.

This is an algebraic diagnostic and an exact model calculation. It does
not replace the analytic estimates needed for arbitrary sparse
majorants. In particular, it is not enough to check the original
configuration if the proof later creates other patterns; all patterns
actually required by that proof must be tested.

\section{Reproducibility, proof dependencies, and integration}
\label{gm:sec:repro}

\subsection{The computational checks}
The supplied \texttt{verify.py} uses Python's standard library and exact
finite-field arithmetic. It does not use numerical ranks, floating-point
character sums, random sampling, or a computer algebra oracle. The
recorded run includes:
\begin{enumerate}[label=(\roman*)]
\item exhaustive enumeration of all $Q^7$ matrices in
\eqref{gm:eq:cube-space} for $Q=2,3,4,5,7$, including extension-field
arithmetic for $\F_4$;
\item all 255 nonempty subsets of the binary three-cube over each of
those fields, checking their ranks and closures;
\item full-grid margin ranks on four grid shapes at every degree up to
the number of axes, and agreement of the margin and symmetric-matrix
representations on selected $2\times3\times2$ patterns;
\item independent direct counting of the seven- and eight-site zero
moments at $n=1$ over $\F_2$ and $\F_3$;
\item formal coefficient expansions for several reduced polynomial
blocks, comparing their full and top-degree coefficient kernels with the
margin kernels; and exhaustive binary checks of the support bound in
three small cases;
\item 28,970 exhaustive weighted-Laplacian frequency checks over
$\F_2$ and $\F_3$, covering four-, six-, and eight-cycles, $K_{2,3}$,
$K_{3,3}$, and a five-edge path. These check the girth exponent, the cycle
formulas, forest equality, and the positive leading coefficients.
\end{enumerate}
The field $\F_4$ is represented by $\F_2[t]/(t^2+t+1)$; the optional
$\F_8$ implementation uses $t^3+t+1$. Prime-field arithmetic uses modular
inverses. The general rank-enumerator theorem is proved by counting in
Section~\ref{gm:sec:enumerator}, not inferred from these finite checks.

\subsection{Dependency structure}
The coefficient-extraction theorem and the support bound are purely
algebraic. The sharp mixing theorem uses them and the elementary
multilinear nonvanishing bound. The closure theorem uses only the
quadratic matrix representation, rank-one rigidity, and finite Fourier
inversion. The graph-girth theorem uses weighted incidence matrices, a two-coordinate
recurrence on a shortest cycle, and an exact reciprocal-weight count.
The cube enumerator additionally uses elementary classification
and counting of nondegenerate bilinear planes, followed by a kernel
incidence count. No proof uses a claimed theorem from the recent
\texttt{openai/math} release as a black box.

\subsection{Suggested integration}
A suitable standalone location is
\begin{quote}\ttfamily\small
Combinatorics/Ramsey/Research/GowersSzemeredi/\\
sparse-grid-moments-quadratic-closure/
\end{quote}
The article's labels have the prefix \texttt{gm:}. The package includes a
source register, theorem-status ledger, build instructions, exact check
outputs, and a checksum manifest. It does not modify existing repository
files, assign itself a source number in the merged report, or alter a
formalization ledger.

A natural Lean development would begin with finite label types and
\texttt{Matrix} over a finite field, proving equality of the margin kernel
and the polynomial-relation kernel. Next would come the one-hot matrix
representation and rank-one rigidity. Those finite algebraic interfaces
can be formalized independently of probability and Fourier analysis.
The supplied Python tests are not a substitute for such proofs.
The graph-girth result provides another independent finite algebraic target:
weighted Laplacian ranks on forests and cycles can be verified before
introducing the moment interpretation.

\section{Further research questions}
\label{gm:sec:questions}

\begin{question}[Closure for larger and irregular patterns]
Can $\cl_2(S)$ be described combinatorially without Gaussian elimination,
for arbitrary partial Cartesian grids? In particular, characterize those
patterns for which every rank-one point in their span is already a sampled
site. Theorem~\ref{gm:thm:closure} shows that precisely this property removes
the first normalized correction.
\end{question}
An efficient structural characterization would turn a quantitative error
question into a local combinatorial test. The seven-site cube demonstrates
that independence alone is insufficient.

\begin{question}[Second corrections beyond the cube]
For a quadratically closed pattern, determine the number of rank-two
matrices in $\mathcal V_S$ from geometric invariants of the site set.
More generally, develop formulas for the coefficients of $R_S$ in terms
of low-rank strata of the span.
\end{question}
Theorem~\ref{gm:thm:girth} already determines the first nonzero order for
all two-axis patterns; its leading coefficient for a general graph remains
an explicit rank count rather than a closed graph formula.
The normalized second coefficient for a general pattern is
\[
 [z^2]R_S(z)-s(Q-1)^2c+\frac{s(s+1)}2(Q-1)^2.
\]
Thus rank-two geometry is the next exact obstacle. It is not determined
by the count of hidden rank-one sites alone.

\begin{question}[Higher-degree extremal Fourier modes]
For $d\geq3$, classify the frequencies off $K_d(S)$ for which
$b_\lambda=\theta_{Q,d}$. Is there a useful analogue of quadratic closure
which controls the first correction after normalization by the one-site
mean?
\end{question}
The proof of the sharp upper bound identifies an equality problem for
multilinear contraction maps. A classification could give higher-degree
versions of Theorem~\ref{gm:thm:closure}, with correction rates that need
not be integral powers of $Q^{-n}$.

\begin{question}[Minimum-support relations with extra axes]
Classify all field-valued relations of support $2^{d+1}$ when $q>d+1$,
and determine their dependence on the characteristic and axis sizes.
\end{question}
The tensor description for $q=d+1$ is complete, but the induction proving
the lower bound does not classify all equality configurations with extra
coordinates. Connections with the trade literature should be checked
before asserting new classifications.

\begin{question}[Beyond $p$-reduced blocks]
What replaces $M_d(S)$ for arbitrary polynomial functions over extension
fields, where Frobenius powers can lower additive degree? Can a matrix of
Frobenius-labelled margins express the exact relation space uniformly in
small characteristic?
\end{question}
The example $z^p$ shows why one cannot simply retain formal degree. A
correct replacement should recover both additive-polynomial relations and
Theorem~\ref{gm:thm:relations} as special cases.

\begin{question}[From block models to general polynomial factors]
Which quantitative high-rank hypotheses imply equidistribution on exactly
$H_d(S)$ for a general family of polynomial factors, with an error governed
by the same relation matrix and with explicit dependence on sparsity?
\end{question}
Theorem~\ref{gm:thm:universality} settles independent copies of one
$p$-reduced block, not arbitrary factors. Extending this to general factors
would connect the model more directly to inverse theorems and relative
counting.

\begin{question}[The torus-chart interface]
Can the exact relation criterion be transferred to the buffered torus
charts in the selected OpenAI source, retaining its stated separation of
rank, cutoff, and approximation costs? Which relation spaces arise under
the actual cloned patterns in that argument?
\end{question}
This requires checking the analytic and geometric hypotheses of that
source; the finite-field calculation alone does not settle it. It is the
most direct route from the present local result back to the motivating
arithmetic-progression program.

\begin{question}[Formal verification and generated certificates]
Can a proof assistant check a finite pattern certificate consisting of
$M_d(S)$, its row reduction, and the hidden-site witnesses, and then invoke
general theorems to certify the corresponding limiting law and quadratic
first correction?
\end{question}
Such a checker would preserve ProveIt's distinction between external
computation and trusted proof. The explicit matrices and identities here
provide a small algebraic starting point before the more demanding
Fourier-probability formalization.

\section*{Conclusion}
\addcontentsline{toc}{section}{Conclusion}
The source condition ``number of axes at most polynomial degree'' is one
visible instance of a more precise finite-field principle: low-order
margin relations determine the limiting configuration law. A dependent
pattern does not become independent merely because the number of blocks
grows. Conversely, many partial grids beyond the axis threshold remain
independent, and their exact status is decidable by a finite matrix.
Quadratic closure adds information that rank alone omits, controlling the
first error remaining after marginal normalization. The graph-girth theorem gives the exact convergence rate throughout the
two-axis regime, and the complete cube
rank enumerator makes this distinction exact in every dimension and
every finite field. These are local results with explicit interfaces and
written proofs, rather than an assertion that the global progression
problem has been settled by this manuscript.

\begin{thebibliography}{9}
\addcontentsline{toc}{section}{References}
\raggedright
\bibitem{gowers}
W.~T. Gowers, \emph{A new proof of Szemer\'edi's theorem},
Geom. Funct. Anal. \textbf{11} (2001), 465--588.
\href{https://doi.org/10.1007/s00039-001-0332-9}{doi:10.1007/s00039-001-0332-9}.

\bibitem{gw}
W.~T. Gowers and J. Wolf, \emph{The true complexity of a system of linear
equations}, Proc. Lond. Math. Soc. (3) \textbf{100} (2010), 155--176.
\href{https://arxiv.org/abs/0711.0185}{arXiv:0711.0185};
\href{https://doi.org/10.1112/plms/pdp019}{doi:10.1112/plms/pdp019}.

\bibitem{cfz}
D. Conlon, J. Fox, and Y. Zhao, \emph{A relative Szemer\'edi theorem},
Geom. Funct. Anal. \textbf{25} (2015), 733--762.
\href{https://arxiv.org/abs/1305.5440}{arXiv:1305.5440};
\href{https://doi.org/10.1007/s00039-015-0324-9}{doi:10.1007/s00039-015-0324-9}.
The inspected PDF is arXiv version 2; in particular Definition 2.8 and
Sections 6.2--6.3.

\bibitem{trades}
E. Ghorbani, S. Kamali, G.~B. Khosrovshahi, and D.~S. Krotov,
\emph{On the volumes and affine types of trades},
Electron. J. Combin. \textbf{27} (2020), no.~1, P1.29.
\href{https://arxiv.org/abs/1810.02296}{arXiv:1810.02296};
\href{https://doi.org/10.37236/8367}{doi:10.37236/8367}.

\bibitem{oa}
OpenAI, \emph{Mathematical manuscript collection}, repository README and
research catalogue, release inspected 7 October 2026.
Commit \texttt{adc7f1241b42e322a6451854ab7e4b4c146bf78a}.
\href{https://github.com/openai/math/tree/adc7f1241b42e322a6451854ab7e4b4c146bf78a}{Pinned repository}.
Catalogue claims are not used as theorems in this article.

\bibitem{oa-grid}
OpenAI, \emph{Quasipolynomial Bounds for Arithmetic Progressions},
manuscript directory dated 23 September 2026, source section
\emph{Detection independent of the current width},
\texttt{build/sections/03-relative-detection.tex},
Lemma with source label \texttt{ap:path:grid}.
\href{https://github.com/openai/math/blob/adc7f1241b42e322a6451854ab7e4b4c146bf78a/preprints/Quasipolynomial-Bounds-for-Arithmetic-Progressions-September-23-2026/build/sections/03-relative-detection.tex}{Pinned source}.
Blob \texttt{ca82af70203a4d9d18371266741b8b7210e658d9}.

\bibitem{proveit}
V. Reshetnikov and contributors, \emph{ProveIt},
\emph{Local Quantitative Refinements of Gowers's Proof of Szemer\'edi's
Theorem}, research overview inspected 7 October 2026.
Repository snapshot \texttt{fb9602e556560a15704b4feb176ba6de08378665}.
\href{https://github.com/VladimirReshetnikov/ProveIt/blob/fb9602e556560a15704b4feb176ba6de08378665/Combinatorics/Ramsey/Research/GowersSzemeredi/local-quantitative-refinements/README.md}{Pinned research overview}.
Research placement does not imply formal verification.
\end{thebibliography}
\end{document}
